% Covariance validity on circulant lattices, and the role of separability.
% Standalone note; compile with pdflatex. Sections are written so they can be
% \input into the supplementary material.
\documentclass[11pt]{article}
\usepackage[margin=1.1in]{geometry}
\usepackage{amsmath, amssymb, bm, booktabs}
\usepackage[round]{natbib}
\newcommand{\Z}{\mathbb{Z}}
\newcommand{\R}{\mathbb{R}}
\newcommand{\dft}{\operatorname{DFT}}
\title{Valid covariance functions on circulant lattices,\\ and what the FFT
  algebra actually requires}
\author{}
\date{2026-09-30}

\begin{document}
\maketitle

\section{Setting and the reviewers' point}

The manuscript builds stationary correlation matrices on a cyclic lattice by
composing a Euclidean correlation function with the great-circle (wrap-around)
distance: with $h_{ik} = \min(|i-k|,\, n-|i-k|)$,
\begin{equation} \label{eq:naive}
  R_{ik} \;=\; \rho(h_{ik}), \qquad
  \rho(h) = \exp\{-(h/\phi)^p\}, \quad p \in (0, 2],
\end{equation}
(eq.~(correlation) in the paper, with ``$p=2$ induces smooth realizations'').
The reviewers correctly note that \eqref{eq:naive} is not a valid construction
in general: a function that is positive definite (PD) on $\R$ does not remain
PD after composition with the geodesic distance on a circle. In the
experiments this was handled by keeping $\phi$ small enough that the matrix
was numerically SPD. This note summarizes the relevant theory, checks which
published experiments were affected (none, it turns out), and proposes
concrete modifications.

\section{When is a circulant correlation matrix valid?}

On the cyclic group $\Z_n$ the answer is exact and cheap to check. A
stationary correlation on $\Z_n$ is a circulant matrix $R$ with base
$r = (r_0, \dots, r_{n-1})$, $r_h = \rho(h)$, and its eigenvalues are the DFT
of the base. The discrete analogue of Bochner's theorem (the Herglotz
theorem; see, e.g., \citealp[Ch.~4]{BrockwellDavis1991}) states:
\begin{quote}
  $R$ is a valid (positive semidefinite) covariance on $\Z_n$ \emph{iff}
  $\dft(r)_k \ge 0$ for all $k = 0, \dots, n-1$.
\end{quote}
Since the code already computes $\dft(r)$ to do FFT algebra, validity can and
should be asserted at construction time for free.

\section{Four provably valid constructions}

\paragraph{(a) Restrict the smoothness exponent: $p \in (0, 1]$.}
\citet{Gneiting2013} (building on \citealp{Schoenberg1942}) characterizes the
isotropic PD functions on spheres $S^d$; the circle $S^1$ contains every
$\Z_n$ as a subset, so PD on $S^1$ implies PD on $\Z_n$ with great-circle
distance. The powered exponential $\exp\{-(h/\phi)^p\}$ is PD on all spheres
\emph{iff} $p \in (0, 1]$, and the geodesic Mat\'ern is PD \emph{iff}
$\nu \le 1/2$. This explains the numerical table below: the paper's $p=1$
choices (image and noise correlations) are valid \emph{at every range
$\phi$}, whereas $p = 1.98$ (blur prior) is not a valid model on the circle
and fails numerically once $\phi$ is large relative to $n$.

\paragraph{(b) Periodize (wrap) the Euclidean covariance.}
For any PD, integrable $C$ on $\R$, the periodization
\begin{equation} \label{eq:wrap}
  C_{\mathrm{per}}(h) \;=\; \sum_{j \in \Z} C(h + jn)
\end{equation}
is PD on $\Z_n$: by Poisson summation, $\dft$ of the wrapped base samples the
(nonnegative) spectral density of $C$. This is the classical device behind
circulant embedding \citep{WoodChan1994, DietrichNewsam1997, Stein1999}. It
preserves the Euclidean family's smoothness, so $p \to 2$ stays available. In
practice two or three terms of \eqref{eq:wrap} suffice ($C$ decays fast), and
the change to \texttt{exp\_cor} is one line:
$r_h = \sum_{j=-J}^{J} \exp\{-(|h + jn|/\phi)^p\}$.

\paragraph{(c) Chordal distance.}
Composing any PD radial function on $\R^2$ with the \emph{chordal} distance
$c(h) = \frac{n}{\pi}\sin(\pi h / n)$ (the straight-line distance between
lattice points placed on a circle of circumference $n$) is PD
\citep{Yadrenko1983, Gneiting2013}. This keeps the closed form
$\exp\{-(c(h)/\phi)^p\}$ valid for all $p \in (0,2]$; chordal and great-circle
distances agree to first order for $h \ll n$, so fitted ranges barely change.

\paragraph{(d) Compact support within the half-period.}
If $C$ is PD on $\R$ and supported on $[0, n/2)$ (e.g.\ Wendland or other
compactly supported families; \citealp{Wendland1995, Gneiting2002}), then
\eqref{eq:wrap} has a single term and the naive construction \eqref{eq:naive}
is \emph{exactly} valid. This is the rigorous version of the ``keep $\phi$
small'' workaround: for the powered exponential the support is infinite, so a
small $\phi$ only makes the negative part of the spectrum exponentially
small, never zero (for $p \approx 2$ the worst eigenvalue behaves like
$-e^{-(n/2\phi)^2}$-ish; see the table). An alternative in the same spirit is
the Gaussian Markov random field route: specify the \emph{precision} directly
on the torus \citep{RueHeld2005}, which is valid by construction and keeps
the circulant structure.

\section{Are the published experiments affected?}

Minimum circulant eigenvalue of the naive base \eqref{eq:naive} and of the
wrapped base \eqref{eq:wrap}:

\begin{center}
\begin{tabular}{rrr rr}
\toprule
$n$ & $p$ & $\phi$ & min eig, naive & min eig, wrapped \\
\midrule
24 & 1.98 & 1.5  & $2.7\cdot 10^{-2}$ & $2.7\cdot 10^{-2}$ \\
24 & 1.98 & 2.0  & $3.6\cdot 10^{-3}$ & $3.6\cdot 10^{-3}$ \\
24 & 1.98 & 6.0  & $-1.5\cdot 10^{-2}$ & $2.6\cdot 10^{-4}$ \\
24 & 1.98 & 10.0 & $-3.7\cdot 10^{-1}$ & $9.1\cdot 10^{-5}$ \\
24 & 1.00 & 10.0 & $3.5\cdot 10^{-2}$ & $5.0\cdot 10^{-2}$ \\
96 & 1.98 & 10.0 & $9.1\cdot 10^{-5}$ & $9.1\cdot 10^{-5}$ \\
\bottomrule
\end{tabular}
\end{center}

The experiments use $p = 1$ with $\phi = 1.5$ for image and noise (valid by
(a) at any range) and $p = 1.98$ with $\phi \in \{1.5, 2\}$ for the blur
prior on $n_v \in \{24, \dots\}$, where the naive minimum eigenvalue is
$\ge 3.6\cdot 10^{-3} > 0$. \textbf{All published experiments therefore used
genuinely SPD matrices}; the issue is that the \emph{recipe} stated in the
paper is not valid in general, and a referee can break it with $p=2$ and a
larger $\phi$.

\section{Suggested modifications to the paper and code}

\begin{enumerate}
\item \textbf{Eq.~(correlation) (Sec.~``Blur prior distribution'').} State the
  validity condition once: ``a stationary correlation on the cyclic lattice
  is valid iff the DFT of its base is nonnegative (Herglotz)''. Then present
  the construction as either (a) $p \in (0,1]$ with great-circle distance
  \citep{Gneiting2013}, or (b) the periodization \eqref{eq:wrap} of the
  Euclidean powered exponential for $p \in (0,2]$, which we recommend since
  the blur prior wants $p \approx 2$ smoothness. Drop ``choosing $p=2$
  induces smooth realizations'' in favour of the wrapped version.
\item \textbf{SM simulation details.} Note that the values used
  ($p = 1.98, \phi \le 2$; $p = 1, \phi = 1.5$) give SPD matrices, so no
  published result changes; for general use the wrapped form should be the
  default.
\item \textbf{Code.} In \texttt{src/utils/model\_utils.py::exp\_cor}
  (circulant branch): replace the base by the $J$-term wrapped sum (b), and
  assert \texttt{np.fft.fft(r).real.min() >= -tol} in
  \texttt{Covariance.create\_exp\_cor\_matrix} so an invalid choice fails
  loudly instead of corrupting a chain.
\end{enumerate}

\section{Separability and the Kronecker product (Reviewer 3, comment 5)}

The reviewer asks why the image covariance is assumed separable
($\bm{R}_c = \bm{R}_h \otimes \bm{R}_v$), suggesting the assumption is
restrictive and perhaps unnecessary. The precise situation, verified against
the implementation:

\paragraph{What the FFT speedups actually need: stationarity, not
separability.} Every matrix in the collapsed likelihood
($A = W \Sigma_c W^{\top} + \Sigma_d$, $Z$, $B$, their inverses, determinants
and mat-vecs) is manipulated only through the 2D DFT of its base. Any
\emph{stationary} covariance on the torus yields a BCCB matrix, which the DFT
diagonalizes whether or not it factorizes as a Kronecker product. The same
holds for the Gibbs full conditionals and for the exact-observation
correction (the gather $A_c C A_c^{\top}$ of a base column works for any
BCCB). So the reviewer is right that the headline $O(n \log n)$ complexity
does not require separability: a nonseparable wrapped 2D Mat\'ern would drop
straight into most of the algebra.

\paragraph{Where separability \emph{is} used.} Three places, all verified in
the code: (i) the model is \emph{specified} through one-dimensional ranges
$(\rho_v, \rho_h)$, i.e.\ separability is a modeling convenience for building
$\bm{R}_c$ from 1D pieces; (ii) the constrained-image factorizations: after
conditioning on a full well column, the conditional covariance of a separable
field factorizes, and the stored $\bm{R}_v^{\star}, \bm{R}_h^{\star}$ (and
the Cholesky factor $L$) exploit this; (iii) the sum-of-squares terms of the
potential gradient use the factorization explicitly
($K = P \bm{R}_h P^{\top}$, $X = \bm{R}_v^{\star} W_0^{\top}$, and the
$\bm{R}_{cv}, \bm{R}_{ch}$ eigenvalue products in the legacy gradient).

\paragraph{Suggested response / rewrite.} Introduce stationarity on the
extended cyclic lattice as the essential assumption (with the validity
condition of Sec.~2) when the model is first presented, and present
separability as an explicit, secondary assumption adopted (i) to elicit the
prior through directional ranges, and (ii) to obtain closed-form
factorizations in the exact-observation algebra and the HMC gradient. It is
fair to concede that the collapsed likelihood itself, and hence Gibbs and the
HMC potential, extend verbatim to nonseparable stationary covariances; the
gradient's sum-of-squares terms would need a (derivable) generalization. In
short: the reviewer's instinct is correct for the core algebra and the paper
overstated the role of the Kronecker product there, but separability is not
vacuous either --- it is what makes the well-constraint and gradient terms
cheap in the current implementation.

\begin{thebibliography}{99}
\bibitem[Brockwell and Davis(1991)]{BrockwellDavis1991}
Brockwell, P.~J. and Davis, R.~A. (1991). \emph{Time Series: Theory and
Methods}, 2nd ed. Springer. (Herglotz theorem, Ch.~4.)
\bibitem[Dietrich and Newsam(1997)]{DietrichNewsam1997}
Dietrich, C.~R. and Newsam, G.~N. (1997). Fast and exact simulation of
stationary Gaussian processes through circulant embedding of the covariance
matrix. \emph{SIAM J. Sci. Comput.} 18(4), 1088--1107.
\bibitem[Gneiting(2002)]{Gneiting2002}
Gneiting, T. (2002). Compactly supported correlation functions.
\emph{J. Multivariate Anal.} 83(2), 493--508.
\bibitem[Gneiting(2013)]{Gneiting2013}
Gneiting, T. (2013). Strictly and non-strictly positive definite functions on
spheres. \emph{Bernoulli} 19(4), 1327--1349.
\bibitem[Rue and Held(2005)]{RueHeld2005}
Rue, H. and Held, L. (2005). \emph{Gaussian Markov Random Fields: Theory and
Applications}. Chapman \& Hall/CRC.
\bibitem[Schoenberg(1942)]{Schoenberg1942}
Schoenberg, I.~J. (1942). Positive definite functions on spheres.
\emph{Duke Math. J.} 9(1), 96--108.
\bibitem[Stein(1999)]{Stein1999}
Stein, M.~L. (1999). \emph{Interpolation of Spatial Data: Some Theory for
Kriging}. Springer.
\bibitem[Wendland(1995)]{Wendland1995}
Wendland, H. (1995). Piecewise polynomial, positive definite and compactly
supported radial functions of minimal degree. \emph{Adv. Comput. Math.} 4,
389--396.
\bibitem[Wood and Chan(1994)]{WoodChan1994}
Wood, A.~T.~A. and Chan, G. (1994). Simulation of stationary Gaussian
processes in $[0,1]^d$. \emph{J. Comput. Graph. Statist.} 3(4), 409--432.
\bibitem[Yadrenko(1983)]{Yadrenko1983}
Yadrenko, M.~I. (1983). \emph{Spectral Theory of Random Fields}.
Optimization Software, New York. (Chordal-distance construction.)
\end{thebibliography}

\end{document}
